% =====================================================================
%  第一部分  数学与物理基础
% =====================================================================
\part{数学与物理基础}

这一部分从最简单的概念开始，一直推到两个控制方程。
如果你已经会偏导数，可以跳过第一章，直接从第二章读起。

% =====================================================================
\section{数学预备知识}

\subsection{一个量可以同时依赖两个变量}

先说一个日常例子。房间里某一点的温度，会随着时间变化，也会随着位置变化。
房间中央和靠窗的位置温度不一样，同一个位置早上和中午也不一样。

为了把这件事写清楚，我们用两个符号：

\begin{itemize}[leftmargin=1.8em,itemsep=2pt]
\item $t$ 表示时间，单位是秒，符号 s。
\item $r$ 表示某个位置到中心的距离，单位是米，符号 m。
\end{itemize}

于是温度就写成一个函数

\begin{equation}
T = T(r,t).
\end{equation}

这个式子的读法是：温度 $T$ 由两个自变量 $r$ 和 $t$ 共同决定。
括号里先写位置，后写时间，这个顺序一旦定下来就不能改。

同样地，水分浓度也是两个自变量的函数，写作

\begin{equation}
C = C(r,t).
\end{equation}

水分浓度的单位是千克每千克，符号 kg/kg，它的确切含义我们在第二章讲。

\subsection{偏导数：把一个变量当常数}

现在的问题是：温度变化有多快？

如果只问时间方向的变化快慢，那就把位置 $r$ 固定住不动，只看 $t$ 变化时 $T$ 怎么变。
这种在一个变量变化、其余变量固定的情况下求出的变化率，叫做偏导数，写作

\begin{equation}
\frac{\partial T}{\partial t}.
\end{equation}

符号 $\partial$ 读作偏，它和普通导数符号 $\mathrm{d}$ 的区别只有一点：
用 $\partial$ 的时候，暗示还有别的自变量存在，只是暂时被固定住了。

计算规则很简单：把另一个变量真的当成一个普通常数去求导就行。举例说明。

假设有一个函数

\begin{equation}
f(x,y) = 3x^2 y + 5y + 7 .
\end{equation}

要求 $\dfrac{\partial f}{\partial x}$，就把 $y$ 当成常数。此时 $3y$ 是常数，
$5y+7$ 也是常数，常数求导为零，于是

\begin{equation}
\frac{\partial f}{\partial x} = 3y \cdot 2x = 6xy .
\end{equation}

要求 $\dfrac{\partial f}{\partial y}$，就把 $x$ 当成常数，于是 $3x^2$ 是常数，

\begin{equation}
\frac{\partial f}{\partial y} = 3x^2 \cdot 1 + 5 = 3x^2 + 5 .
\end{equation}

回到温度。$\dfrac{\partial T}{\partial t}$ 的含义是：
站在药材的某一个固定位置上，看温度随时间涨得多快，单位是摄氏度每秒。

$\dfrac{\partial T}{\partial r}$ 的含义是：
在某个固定时刻，沿着半径方向走一小段，看温度变化多少，单位是摄氏度每米。
这个量叫温度梯度，它是热量流动的直接原因。

\subsection{二阶偏导数}

偏导数本身还是一个函数，所以可以再对它求一次偏导数。
对同一个变量求两次，得到

\begin{equation}
\frac{\partial^2 T}{\partial r^2}
= \frac{\partial}{\partial r}\left(\frac{\partial T}{\partial r}\right).
\end{equation}

这个量的物理含义是温度梯度沿半径方向的变化快慢。
如果 $\dfrac{\partial^2 T}{\partial r^2}$ 大于零，说明温度曲线是向上弯的，
梯度在往外走的过程中越来越大。

\subsection{圆柱坐标与那一项多出来的东西}

平面上，一个点到原点的距离记作 $r$，它和直角坐标 $x$、$y$ 的关系是

\begin{equation}
r = \sqrt{x^2+y^2}.
\end{equation}

药材是圆柱形的，它的截面是一个圆，温度只和到圆心的距离 $r$ 有关，
和转过的角度无关，这种情形叫做轴对称。

在直角坐标里，温度对位置的全部二阶变化写成一个算子，叫拉普拉斯算子，

\begin{equation}
\nabla^2 T = \frac{\partial^2 T}{\partial x^2}+\frac{\partial^2 T}{\partial y^2}.
\end{equation}

换到圆柱坐标以后，这个算子变成

\begin{equation}
\nabla^2 T = \frac{1}{r}\frac{\partial}{\partial r}\left(r\frac{\partial T}{\partial r}\right).
\label{eq:lap}
\end{equation}

这个式子里的 $1/r$ 和那个多出来的 $r$ 是从哪里来的，值得说清楚，
因为它正是后面第五章要大动干戈的原因。

推导过程如下。圆柱坐标下，面积微元是

\begin{equation}
\mathrm{d}A = r\,\mathrm{d}r\,\mathrm{d}\theta .
\end{equation}

径向的热流密度乘以周长 $2\pi r$ 才是总热流，所以圆周长里带一个 $r$。
把通量沿径向的净变化除以面积 $2\pi r\,\mathrm{d}r$ 时，
分子分母各有一个 $r$ 约不掉，最后就剩下这个形式。第二章会把这个过程完整写出来。

把式 \eqref{eq:lap} 展开，得到另一种写法：

\begin{equation}
\frac{1}{r}\frac{\partial}{\partial r}\left(r\frac{\partial T}{\partial r}\right)
= \frac{1}{r}\left[\frac{\partial T}{\partial r} + r\frac{\partial^2 T}{\partial r^2}\right]
= \frac{\partial^2 T}{\partial r^2} + \frac{1}{r}\frac{\partial T}{\partial r}.
\label{eq:lap-expand}
\end{equation}

中间那一步用到了乘积求导法则。设 $u=r$，$v=\dfrac{\partial T}{\partial r}$，
则 $(uv)' = u'v + uv'$，也就是

\begin{equation}
\frac{\partial}{\partial r}\left(r\frac{\partial T}{\partial r}\right)
= 1\cdot\frac{\partial T}{\partial r} + r\cdot\frac{\partial^2 T}{\partial r^2}.
\end{equation}

再除以 $r$ 就得到式 \eqref{eq:lap-expand}。

\keypoint
式 \eqref{eq:lap-expand} 里出现了 $1/r$。
当 $r$ 趋向于零的时候，这一项看起来会变成无穷大。
这是圆柱坐标带来的麻烦，第五章会用一个巧妙的换元把它消掉。

\subsection{微分方程、阶、线性与非线性}

一个方程里如果含有未知函数的导数，就叫微分方程。
我们要求解的 $T$ 和 $C$ 都是未知函数，所以下面出现的都是微分方程。

按最高阶导数的阶数命名。只出现一阶导数叫一阶方程，
出现二阶导数叫二阶方程。本题的两个方程都是二阶的。

按未知函数和它的导数是否以一次方的形式出现，分为线性与非线性。

如果方程里的系数只依赖自变量，比如 $r$ 和 $t$，那么方程是线性的。
线性方程有一个很好的性质：两个解加起来还是解。
如果系数依赖未知函数本身，比如导热系数依赖温度，那就成了非线性方程，
线性方程的许多便利就不能用了。

本题中，问题一的所有物性都是常数，所以方程是线性的，
而且是两个互不影响的线性方程。问题二到问题四的物性依赖 $C$ 和 $T$，
方程变成非线性，两个方程还互相耦合，必须一起求解。

\subsection{初值条件与边界条件}

一个二阶微分方程要有唯一确定的解，需要补充条件。

对时间方向，因为方程里出现 $\dfrac{\partial T}{\partial t}$，
需要一个起始时刻的状态，叫做初值条件，也叫初始条件：

\begin{equation}
T(r,0)=28\,^\circ\mathrm{C}, \qquad C(r,0)=2.55\ \mathrm{kg/kg},
\qquad 0\le r\le R_0 .
\end{equation}

读法是：在时间等于零的时刻，从圆心到表面每一个位置的温度和水分浓度都已知。
题目给出的是这两组数。

对空间方向，因为方程里出现 $\dfrac{\partial^2 T}{\partial r^2}$，
需要在空间区域的两个端点各给一个条件。本题的空间区域是从圆心 $r=0$
到表面 $r=R_0$，所以要在 $r=0$ 和 $r=R_0$ 各给一个条件，这叫边界条件。

\subsection{三类边界条件}

边界条件通常分三类，理解这三类对读懂本题很关键。

\textbf{第一类，给定边界上的值。} 例如直接规定表面温度恒为 $80$ 摄氏度，
写作 $T(R_0,t)=80$。这类条件叫狄利克雷条件。

\textbf{第二类，给定边界上的通量。} 例如规定表面每平方米每秒流入多少焦耳的热量，
写作 $-k\dfrac{\partial T}{\partial r}\Big|_{r=R_0}=q_0$。
这里出现的是导数而不是函数值。这类条件叫诺伊曼条件。

\textbf{第三类，给定边界值与外界的关系。} 例如表面通过空气对流换热，
换热快慢取决于表面和空气的温差，写作

\begin{equation}
-k\left.\frac{\partial T}{\partial r}\right|_{r=R_0} = h\left(T(R_0,t)-T_{\mathrm{air}}(t)\right).
\end{equation}

这类条件叫罗宾条件，本题的表面条件正是这一种，
因为热风烘干是靠空气流动把热量带到药材表面的。

\keypoint
初值条件管时间方向的起点，边界条件管空间方向的两端。
两类条件都齐全，方程的解才唯一。

% =====================================================================
\section{从守恒律推出控制方程}

这一章是整份文档的核心。我们从一个薄圆筒壳出发，
用中学物理里的能量守恒和质量守恒，一步步推出两个控制方程。
每一步的代数运算都会写全。

\subsection{微元法的总思路}

直接对整个圆柱列方程很难，因为圆柱内部各处的温度不一样。
办法是取一个很小的区域，让这个小区域内部可以近似认为温度是均匀的，
对这个小区域列方程，然后再让它缩小到一点，得到只描述该点的方程。

这个小区域叫微元，也叫控制体。取法有很多种，本题取一个薄圆筒壳，
因为它正好配合圆柱的轴对称性质。

\subsection{傅里叶定律}

热量总是从高温处流向低温处，流动的快慢与温度梯度成正比。
这个经验规律叫傅里叶定律，一维形式写作

\begin{equation}
q_r = -k\frac{\partial T}{\partial r}.
\label{eq:fourier}
\end{equation}

其中 $q_r$ 是沿半径方向的热流密度，单位是瓦特每平方米，符号 $\mathrm{W/m^2}$，
它表示单位时间穿过单位面积的热量。
$k$ 是导热系数，单位是瓦特每米每开尔文，符号 $\mathrm{W/(m\cdot K)}$，
它衡量材料导热能力的强弱。

式 \eqref{eq:fourier} 里的负号表示方向：
如果温度沿半径向外降低，也就是 $\dfrac{\partial T}{\partial r}<0$，
那么 $q_r>0$，热量朝外流，方向和温度降低的方向一致。
如果温度沿半径向外升高，热量就朝内流。

\subsection{圆柱薄壳的几何量}

现在取微元。半径从 $r$ 到 $r+\mathrm{d}r$ 之间的一层薄壳，
高度取药材的长度 $L$。

这个薄壳的体积要算清楚，因为后面要用它表示内能。

把薄壳看成一张卷起来的纸，展开以后它的长度是圆周长 $2\pi r$，
高度是 $L$，厚度是 $\mathrm{d}r$。所以体积是

\begin{equation}
\mathrm{d}V = 2\pi r L\,\mathrm{d}r .
\label{eq:dV}
\end{equation}

这里出现的 $2\pi r$ 就是圆周长，那个 $r$ 来自周长公式。
请特别记住这一点，它是圆柱坐标里所有 $r$ 因子的来源。

薄壳的质量等于体积乘密度，

\begin{equation}
\mathrm{d}m = \rho\,\mathrm{d}V = \rho\,2\pi r L\,\mathrm{d}r ,
\label{eq:dm}
\end{equation}

其中 $\rho$ 是密度，单位是千克每立方米，符号 $\mathrm{kg/m^3}$。

\subsection{能量守恒，一步不跳}

对这个薄壳列能量守恒。表述是：

\begin{equation}
\text{薄壳内能的增加速率} = \text{单位时间净流入薄壳的热量}.
\end{equation}

先算左边。内能可以用比热容表示。比热容 $c_p$ 的含义是：
让一千克物质升高一摄氏度需要吸收多少焦耳的热量，
单位是焦耳每千克每开尔文，符号 $\mathrm{J/(kg\cdot K)}$。

那么，质量为 $\mathrm{d}m$ 的物质，温度变化率为 $\dfrac{\partial T}{\partial t}$，
内能增加速率就是

\begin{equation}
\text{内能增加速率} = c_p\,\mathrm{d}m\,\frac{\partial T}{\partial t}.
\label{eq:lhs}
\end{equation}

把式 \eqref{eq:dm} 代进去，得到

\begin{equation}
\text{左边} = \rho c_p\,2\pi r L\,\mathrm{d}r\,\frac{\partial T}{\partial t}.
\label{eq:lhs2}
\end{equation}

再算右边。热量从两个面进出：内表面 $r$ 处和外表面 $r+\mathrm{d}r$ 处。

约定朝外为正方向。那么在 $r$ 处流入薄壳的热量，等于该处朝外的热流取负号：
$-q_r(r)\cdot 2\pi r L$。这里 $2\pi r L$ 是内表面的面积。

在 $r+\mathrm{d}r$ 处流出薄壳的热量是 $+q_r(r+\mathrm{d}r)\cdot 2\pi(r+\mathrm{d}r)L$。

两者相加就是净流入：

\begin{equation}
\text{右边} = -q_r(r)\cdot 2\pi r L + q_r(r+\mathrm{d}r)\cdot 2\pi(r+\mathrm{d}r)L .
\label{eq:rhs}
\end{equation}

这个式子可以改写成更紧凑的形式。注意右边正是函数 $q_r(r)\cdot 2\pi r L$
在区间两端点的差，也就是

\begin{equation}
\text{右边} = \Big[q_r\cdot 2\pi r L\Big]_{r}^{r+\mathrm{d}r}.
\label{eq:rhs2}
\end{equation}

让左右两边相等，得到

\begin{equation}
\rho c_p\,2\pi r L\,\mathrm{d}r\,\frac{\partial T}{\partial t}
= \Big[q_r\cdot 2\pi r L\Big]_{r}^{r+\mathrm{d}r}.
\label{eq:balance1}
\end{equation}

\subsection{从差分到微分}

式 \eqref{eq:balance1} 的右边是一个差值，还不能直接用。
回忆微积分里的一个基本事实：函数在区间上的增量，
除以区间长度，在区间长度趋向于零时等于导数，即

\begin{equation}
\lim_{\mathrm{d}r\to0}\frac{f(r+\mathrm{d}r)-f(r)}{\mathrm{d}r} = \frac{\mathrm{d}f}{\mathrm{d}r}.
\end{equation}

所以，把式 \eqref{eq:balance1} 两边同时除以 $\mathrm{d}r$，得到

\begin{equation}
\rho c_p\,2\pi r L\,\frac{\partial T}{\partial t}
= \frac{\Big[q_r\cdot 2\pi rL\Big]_{r}^{r+\mathrm{d}r}}{\mathrm{d}r}.
\end{equation}

再让 $\mathrm{d}r$ 趋向于零，右边就变成导数

\begin{equation}
\rho c_p\,2\pi r L\,\frac{\partial T}{\partial t}
= \frac{\partial}{\partial r}\left(q_r\cdot 2\pi rL\right).
\end{equation}

式子两边都有因子 $2\pi L$，可以约掉。约掉以后

\begin{equation}
\rho c_p\,r\,\frac{\partial T}{\partial t}
= \frac{\partial}{\partial r}\left(q_r\,r\right).
\label{eq:balance2}
\end{equation}

现在把傅里叶定律 \eqref{eq:fourier} 代入，即 $q_r=-k\dfrac{\partial T}{\partial r}$，

\begin{equation}
\rho c_p\,r\,\frac{\partial T}{\partial t}
= \frac{\partial}{\partial r}\left(-k\,r\,\frac{\partial T}{\partial r}\right).
\end{equation}

把右边的负号移出导数符号外面，

\begin{equation}
\rho c_p\,r\,\frac{\partial T}{\partial t}
= -\frac{\partial}{\partial r}\left(k\,r\,\frac{\partial T}{\partial r}\right).
\end{equation}

两边同乘 $-1$，并注意 $k$ 在问题一中是常数，可以移到导数外面，

\begin{equation}
-\rho c_p\,r\,\frac{\partial T}{\partial t}
= \frac{\partial}{\partial r}\left(k\,r\,\frac{\partial T}{\partial r}\right).
\end{equation}

这里出现了符号不一致的问题。我们换个做法：不把负号移出去，
而是两边同除以 $r$，得到

\begin{equation}
\rho c_p\,\frac{\partial T}{\partial t}
= \frac{1}{r}\frac{\partial}{\partial r}\left(r\,k\,\frac{\partial T}{\partial r}\right).
\end{equation}

这个式子形式漂亮，而且物理意义清楚：左边是单位体积的内能变化率，
右边是净流入的热量。最后把它整理成标准写法

\begin{equation}
\boxed{\;
\rho(C)\,c_p(C)\,\frac{\partial T}{\partial t}
= \frac{1}{r}\frac{\partial}{\partial r}\left(r\,k(C)\,\frac{\partial T}{\partial r}\right).
\;}
\label{eq:energy}
\end{equation}

这就是第一个控制方程。

我们把系数写成 $\rho(C)$、$c_p(C)$、$k(C)$ 是为了提醒：
在问题二到问题四里，它们都随水分浓度 $C$ 变化，不是常数。
在问题一里它们都是常数，方程是线性的。

\why{为什么右边的 $r$ 不能约掉}
左边分子分母各有一个 $r$ 看起来可以约，实际上不行。
左边的 $r$ 来自体积公式里的圆周长，右边的 $r$ 在导数符号里面，
它随 $r$ 变化，不是常数因子，不能拿到导数外面去。
两者只是形式上一样，物理来源不同。

\subsection{菲克定律与水分守恒}

水分的迁移和热量类似，也是由浓度差驱动，规律叫菲克定律。一维形式写作

\begin{equation}
j_r = -D\frac{\partial C}{\partial r},
\end{equation}

其中 $j_r$ 是水分通量，$D$ 是水分扩散系数，单位是平方米每秒，符号 $\mathrm{m^2/s}$。
负号的含义和傅里叶定律一样：水分从浓度高处流向浓度低处。

但是这里有一个容易出错的地方。水分是存在药材内部的，
药材本身在干燥过程中还会收缩，所以不能简单地把水分通量除以面积就完事。
必须盯住干物质，因为干物质是不迁移的。

正确的做法是，用干基含水率 $C$ 作为未知量，
并且注意到控制体里的干物质质量不随时间变化。
第二章下一小节详细说这件事。

\subsection{干基含水率是什么}

日常生活中说含水率，通常指湿基含水率，定义是

\begin{equation}
\text{湿基含水率} = \frac{\text{水的质量}}{\text{水的质量}+\text{干物质的质量}} .
\end{equation}

题目用的是干基含水率，定义是

\begin{equation}
C = \frac{\text{水的质量}}{\text{干物质的质量}} .
\label{eq:drybasis}
\end{equation}

例如题目给初始值 $C_0=2.55$ kg/kg，意思是每 $1$ 千克干物质里含有 $2.55$ 千克水。

\why{为什么用干基而不是湿基}
因为干燥过程中水分不断减少，用湿基的话分母一直在变，
写守恒方程时会出现额外的项。用干基的话，分母是干物质质量，
而干物质不迁移也不增减，是一个固定的参考量。
这样写出来的方程最干净，处理收缩问题也最方便。
这是一个非常重要但常被忽略的建模技巧。

\subsection{干物质不迁移这个事实}

现在把控制体从按空间位置取，改成按干物质取。
具体说，我们跟踪同一批干物质，看它内部的水分怎么变。

这样取控制体的好处是：控制体自身的质量不变，
水分的进出只能通过控制体的两个侧面，不会因为控制体自己收缩而凭空多出或少掉水。

单位湿体积内的干物质质量记作

\begin{equation}
\rho_d = \frac{\rho}{1+C}.
\end{equation}

这个式子怎么来的？体积内总质量是 $\rho$，其中干物质占的比例是
$\dfrac{1}{1+C}$。因为按式 \eqref{eq:drybasis}，每 $1$ 千克干物质对应 $C$ 千克水，
所以总质量是干物质的 $1+C$ 倍。

\subsection{水分守恒方程}

仿照传热的做法。控制体内水分的变化率等于净流入的水分。

水分通量按干物质计，写成

\begin{equation}
j_r = -\rho_d D\frac{\partial C}{\partial r}.
\end{equation}

单位时间流入控制体的水分是 $j_r$ 乘以侧面积 $2\pi rL$。
控制体内的水分质量是 $C$ 乘以干物质质量，而干物质质量不变，
所以水分质量变化率是

\begin{equation}
\text{干物质质量}\times\frac{\partial C}{\partial t}.
\end{equation}

于是守恒式是

\begin{equation}
\rho_d\,2\pi rL\,\mathrm{d}r\,\frac{\partial C}{\partial t}
= \Big[j_r\cdot 2\pi rL\Big]_{r}^{r+\mathrm{d}r}.
\end{equation}

和前面完全一样的步骤：除以 $\mathrm{d}r$，让 $\mathrm{d}r$ 趋于零，
两边约掉 $2\pi L$ 和 $\rho_d$，得到

\begin{equation}
r\,\frac{\partial C}{\partial t}
= \frac{\partial}{\partial r}\left(r\,D\,\frac{\partial C}{\partial r}\right),
\end{equation}

再除以 $r$，

\begin{equation}
\frac{\partial C}{\partial t}
= \frac{1}{r}\frac{\partial}{\partial r}\left(r\,D\,\frac{\partial C}{\partial r}\right).
\end{equation}

最后写成标准形式

\begin{equation}
\boxed{\;
\frac{\partial C}{\partial t}
= \frac{1}{r}\frac{\partial}{\partial r}\left(r\,D(C,T)\,\frac{\partial C}{\partial r}\right).
\;}
\label{eq:mass}
\end{equation}

这就是第二个控制方程。

\why{两个方程的左边为什么不同}
传热方程的左边是 $\rho c_p\dfrac{\partial T}{\partial t}$，
水分方程的左边只有 $\dfrac{\partial C}{\partial t}$。
原因在于未知量的选取方式不同。
温度是单位质量的能量指标，要乘上密度和比热容才变成单位体积的能量。
而 $C$ 已经是每千克干物质的水量，本身就是一个比值，
控制体又跟着干物质走，所以左边干净。
这两点如果搞混，方程就会写错。

\subsection{无量纲分析：先判断哪边快}

在动手算之前，先做一件事：估算两个过程的快慢，看哪个先完成。

定义一个时间尺度，表示某个过程走完全程大概需要多久。
传热的特征时间取

\begin{equation}
\tau_T = \frac{R_0^2}{\alpha}, \qquad \alpha = \frac{k}{\rho c_p},
\end{equation}

其中 $\alpha$ 叫热扩散率，单位是平方米每秒。
传质的特征时间取

\begin{equation}
\tau_D = \frac{R_0^2}{D}.
\end{equation}

代入题目数据，问题一里

\begin{equation}
\alpha = \frac{0.36}{820\times2600} = 1.6886\times10^{-7}\ \mathrm{m^2/s},
\end{equation}

\begin{equation}
\tau_T = \frac{0.02^2}{1.6886\times10^{-7}} = 2369\ \mathrm{s},
\end{equation}

而水分扩散系数在 $C=2.55$ 时等于

\begin{equation}
D = 7\times10^{-9}\times e^{-0.89/2.55} = 7\times10^{-9}\times0.7053
= 4.9377\times10^{-9}\ \mathrm{m^2/s},
\end{equation}

\begin{equation}
\tau_D = \frac{0.02^2}{4.9377\times10^{-9}} = 8.10\times10^4\ \mathrm{s}.
\end{equation}

两者相差三十多倍。结论很清楚：温度很快就能平衡，水分则慢得多。
这个判断在第六章解释结果时会反复用到。

\keypoint
传热特征时间约 $2369$ 秒，传质特征时间约 $81000$ 秒。
所以题目里说的前 $30$ 分钟，温度明显在上升，水分却几乎没有动。

% =====================================================================
\section{边界条件与初始条件}

方程本身只描述了内部各点之间的相互关系，还需要补充端点与起始时刻的信息，
才能得到唯一的解。这一章把四个条件一个一个推导出来。

\subsection{圆心处的条件：对称性}

在圆心 $r=0$ 处，温度不可能是无穷大，也不能有尖角。
理由是对称性：圆心左边一点点和右边一点点到圆心的距离都是 $r$，
由轴对称假设它们的温度必须相同。也就是说温度是 $r$ 的偶函数，
偶函数在原点处的导数为零。所以

\begin{equation}
\left.\frac{\partial T}{\partial r}\right|_{r=0}=0,
\qquad
\left.\frac{\partial C}{\partial r}\right|_{r=0}=0 .
\label{eq:sym}
\end{equation}

\why{这个条件为什么不能省}
如果不用它，数学上会出现无穷多个解，其中大多数在物理上荒谬，
比如温度在圆心附近无限升高。
另外，它还保证了第一章里提到的 $1/r$ 那一项在原点不会真的发散，
因为分子 $\dfrac{\partial T}{\partial r}$ 同时趋向于零。

\subsection{表面处的条件：牛顿冷却定律}

药材表面与热空气接触，热量通过空气流动带过来。这个过程叫对流换热。
牛顿冷却定律说，单位面积上单位时间传入的热量与温差成正比，

\begin{equation}
q_{\text{传入}} = h\left(T_{\mathrm{air}}(t)-T(R_0,t)\right),
\label{eq:newton}
\end{equation}

其中 $h$ 叫对流换热系数，单位是瓦特每平方米每开尔文，
符号 $\mathrm{W/(m^2\cdot K)}$。题目给出 $h=25$。

把式 \eqref{eq:newton} 转化成我们方程里的量。
朝外为正的热流密度是 $q_r=-k\dfrac{\partial T}{\partial r}$。
朝内传入的热量等于朝外热流的负值，所以

\begin{equation}
-k\left.\frac{\partial T}{\partial r}\right|_{r=R_0}
= h\left(T(R_0,t)-T_{\mathrm{air}}(t)\right).
\label{eq:robinT}
\end{equation}

左边是导热给出的通量，右边是对流给出的通量，两者必须相等，
否则表面处的能量就不守恒。

同样的推理用在水分上。水分由空气带走，传质系数记作 $h_m$，
单位是米每秒，符号 $\mathrm{m/s}$，题目给出 $h_m=8\times10^{-7}$。

朝外为正的水分通量是 $-D\dfrac{\partial C}{\partial r}$，于是

\begin{equation}
-D\left.\frac{\partial C}{\partial r}\right|_{r=R_0}
= h_m\left(C(R_0,t)-C_{\mathrm{air}}(t)\right).
\label{eq:robinC}
\end{equation}

\why{为什么左右两边的单位能对上}
以温度为例。左边是导热系数乘以温度梯度，
单位是 $\mathrm{W/(m\cdot K)}\times\mathrm{K/m}=\mathrm{W/m^2}$。
右边是换热系数乘以温差，
单位是 $\mathrm{W/(m^2\cdot K)}\times\mathrm{K}=\mathrm{W/m^2}$。
两边都是热流密度，单位一致。这是一次必要的量纲检查，
如果对不上就说明公式写错了。

\subsection{初始条件}

题目给出烘干开始时药材温度 $28$ 摄氏度，干基含水率 $2.55$ 千克每千克，
且药材内部各处都相同。于是

\begin{equation}
T(r,0)=28\,^\circ\mathrm{C},\qquad C(r,0)=2.55\ \mathrm{kg/kg},
\qquad 0\le r\le R_0 .
\label{eq:ic}
\end{equation}

这里的 $0\le r\le R_0$ 表示整个径向范围都要满足，
不是只在表面或者只在圆心。

\subsection{烘房环境数据的处理}

边界条件 \eqref{eq:robinT} 和 \eqref{eq:robinC} 里出现了两个量，
烘房温度 $T_{\mathrm{air}}(t)$ 和烘房水分浓度 $C_{\mathrm{air}}(t)$。
题目没有给出公式，而是给了附件一，里面是一张表：

\begin{center}
\begin{tabular}{@{}rrr@{}}
\toprule
时间 / s & 温度 / $^\circ$C & 水分浓度 / (kg/kg) \\
\midrule
0    & 28.000 & 0.01963 \\
60   & 28.528 & 0.02002 \\
120  & 29.192 & 0.02053 \\
$\vdots$ & $\vdots$ & $\vdots$ \\
14400 & 50.165 & 0.04986 \\
\bottomrule
\end{tabular}
\end{center}

一共 $241$ 行，每 $60$ 秒一行，时间从 $0$ 到 $14400$ 秒。

问题来了：方程里的 $T_{\mathrm{air}}(t)$ 需要任意时刻的值，
而表格只给了 $60$ 秒一个点。中间那些时刻怎么办？

有两种做法。

第一种做法是直接插值。两个相邻数据点之间连一条直线，
任意时刻的值取直线上的对应点。这种做法最简单，但把测量噪声也带进来了。

第二种做法是先识别数据的规律，用一条光滑曲线去代表它，
再用这条曲线作为边界条件。这种做法叫数据拟合。

从表里能看出明显的规律：温度从一个较低的值开始，
快速上升，然后逐渐变慢，最后稳定在一个固定值附近。
这种形状叫一阶惯性上升。它的数学形式是

\begin{equation}
T_{\mathrm{air}}(t)=T_{\mathrm{set}}-\left(T_{\mathrm{set}}-28\right)e^{-t/\tau_T}.
\label{eq:fitted-T}
\end{equation}

这个式子有三个参数，逐个解释。

$T_{\mathrm{set}}$ 是最终稳定的温度。当 $t$ 很大时，$e^{-t/\tau_T}$ 趋向于零，
式子只剩下 $T_{\mathrm{set}}$，所以它代表终值。

$28$ 是初始温度，直接取表格第一行的实测值，不参与拟合。

$\tau_T$ 叫时间常数，它控制上升的快慢，单位是秒。
当 $t=\tau_T$ 时，指数项等于 $e^{-1}=0.368$，
也就是说温度已经走完了全程的 $63.2\%$。

含湿量用同样的形式，

\begin{equation}
C_{\mathrm{air}}(t)=C_{\mathrm{set}}-\left(C_{\mathrm{set}}-0.01963\right)e^{-t/\tau_C}.
\label{eq:fitted-C}
\end{equation}

\subsection{最小二乘拟合是怎么做的}

现在要确定 $T_{\mathrm{set}}$ 和 $\tau_T$ 这两个数，
使得曲线尽量贴近表格里的 $241$ 个实测点。

贴近的程度用残差平方和衡量。对第 $i$ 个数据点，记实测温度是 $y_i$，
时间是 $t_i$，模型给出的值是 $T_{\mathrm{air}}(t_i)$，则残差是

\begin{equation}
r_i = y_i - T_{\mathrm{air}}(t_i).
\end{equation}

残差平方和是

\begin{equation}
S(T_{\mathrm{set}},\tau_T)=\sum_{i=1}^{241} r_i^2
= \sum_{i=1}^{241}\left[y_i-\left(T_{\mathrm{set}}-\left(T_{\mathrm{set}}-28\right)e^{-t_i/\tau_T}\right)\right]^2 .
\label{eq:lsq}
\end{equation}

我们要找让 $S$ 最小的那一组参数。这就是最小二乘法的含义：
让所有点的偏差平方加起来最小。

\why{为什么用平方而不是绝对值}
用平方有两个好处。一是平方函数处处可导，可以用求导的办法找最小值；
绝对值在零点不可导，处理起来麻烦。二是平方会让大偏差受到更重的惩罚，
拟合结果不会为了照顾某个点而牺牲整体。

\why{为什么把 $S$ 对参数求导等于零就能找到最小值}
在最小值点，函数沿任何方向的切线都是水平的，也就是偏导数为零。
所以令

\begin{equation}
\frac{\partial S}{\partial T_{\mathrm{set}}}=0,\qquad
\frac{\partial S}{\partial \tau_T}=0,
\end{equation}

解这个方程组就得到最优参数。因为式 \eqref{eq:lsq} 里 $\tau_T$ 出现在指数上，
这两个方程无法用代数方法解出来，只能用迭代的办法逐步逼近。
这种算法叫 Levenberg–Marquardt 算法，是数值优化的常用方法。
实际计算时我们调用现成的程序，不需要自己写。

\subsection{拟合结果}

对附件一的数据做上述拟合，得到

\begin{center}
\begin{tabular}{@{}lcccc@{}}
\toprule
参数 & 拟合值 & 标准误 & 均方根误差 & 决定系数 \\
\midrule
$T_{\mathrm{set}}$ & $50.2121\ ^\circ$C & $0.0303$ & \multirow{2}{*}{$0.336\ ^\circ$C} & \multirow{2}{*}{$0.9957$} \\
$\tau_T$ & $1877.49$ s & $12.74$ & & \\
$C_{\mathrm{set}}$ & $0.050908$ kg/kg & $9.3\times10^{-5}$ & \multirow{2}{*}{$8.1\times10^{-4}$} & \multirow{2}{*}{$0.9901$} \\
$\tau_C$ & $2776.34$ s & $31.71$ & & \\
\bottomrule
\end{tabular}
\end{center}

表里几个量要解释一下。

标准误衡量这个参数估得准不准。$T_{\mathrm{set}}$ 的标准误是 $0.0303$ 摄氏度，
意思是这个值大概有正负 $0.03$ 度的不确定度。

均方根误差是残差平方和除以点数再开方，

\begin{equation}
\mathrm{RMSE}=\sqrt{\frac{1}{n}\sum_{i=1}^{n}r_i^2},
\end{equation}

它表示平均每个点偏离曲线多少。温度的平均偏差是 $0.336$ 摄氏度，
这个量级和测量仪表的噪声相当，说明拟合是合理的。

决定系数 $R^2$ 衡量曲线解释了数据的多大比例，定义是

\begin{equation}
R^2 = 1-\frac{\sum_{i=1}^{n}r_i^2}{\sum_{i=1}^{n}\left(y_i-\bar{y}\right)^2},
\end{equation}

其中 $\bar{y}$ 是实测值的平均。$R^2$ 越接近 $1$ 越好。
温度是 $0.9957$，含湿量是 $0.9901$，都在 $0.99$ 以上，说明拟合质量可以接受。

\why{计算时为什么把参数取成 50.212 而不是 50.2121}
拟合给出的完整数值是 $50.2121$，计算时取了 $50.212$。
这样做的影响在表面温度上大约是 $6\times10^{-5}$ 摄氏度，
而拟合本身的标准误是 $0.030$ 摄氏度，前者比后者小三个数量级。
换句话说，纠结最后一位小数是没有意义的，
因为标定数据本身就没那么准。但为了结果能与他人核对，
文档中明确写出使用了哪一组数值。

\subsection{工艺分段}

观察附件一可以发现，温度在 $t$ 大约等于 $4$ 小时，也就是 $14400$ 秒时，
已经升到接近终值并稳定下来。题目也说整个烘干过程分为预热平衡和恒温干燥两个阶段。

所以从 $14400$ 秒往后，烘房环境就当作恒定的：

\begin{equation}
(T_{\mathrm{air}},C_{\mathrm{air}})=
\begin{cases}
\text{按式 \eqref{eq:fitted-T} 和 \eqref{eq:fitted-C} 变化}, & 0\le t\le 14400\ \mathrm{s},\\
\left(50.212\,^\circ\mathrm{C},\ 0.050908\ \mathrm{kg/kg}\right), & t>14400\ \mathrm{s}.
\end{cases}
\label{eq:stage}
\end{equation}

\why{恒温阶段为什么含湿量也不变}
烘房在持续排湿。如果不排湿，药材蒸发出来的水分会留在烘房里，
含湿量会不断上升，最后趋于饱和，干燥就无法继续。
实际生产中排湿速率会把含湿量稳定在一个较低的值上，
所以当作常数处理是合理的。

\subsection{第一部分小结}

到这里，模型已经完整了。把所有式子集中写在一起：

\begin{center}
\begin{tabular}{@{}ll@{}}
\toprule
类别 & 式子 \\
\midrule
传热方程 & $\rho c_p\dfrac{\partial T}{\partial t}=\dfrac{1}{r}\dfrac{\partial}{\partial r}\left(rk\dfrac{\partial T}{\partial r}\right)$ \\
传质方程 & $\dfrac{\partial C}{\partial t}=\dfrac{1}{r}\dfrac{\partial}{\partial r}\left(rD\dfrac{\partial C}{\partial r}\right)$ \\
圆心条件 & $\dfrac{\partial T}{\partial r}=\dfrac{\partial C}{\partial r}=0$，当 $r=0$ \\
表面条件 & $-k\dfrac{\partial T}{\partial r}=h\left(T-T_{\mathrm{air}}\right)$，当 $r=R_0$ \\
表面条件 & $-D\dfrac{\partial C}{\partial r}=h_m\left(C-C_{\mathrm{air}}\right)$，当 $r=R_0$ \\
初始条件 & $T=28$，$C=2.55$，当 $t=0$ \\
\bottomrule
\end{tabular}
\end{center}

问题二和问题三用同一组方程，只是物性换成附录三给出的关系式。
问题四还要加上半径随时间变化这一条，那部分在第二部分第九章讲，本章没有展开。

\keypoint
第一部分讲完了从物理守恒到两个偏微分方程的全过程。
核心有三点。第一，微元法加守恒律给出方程。
第二，干基含水率让水分方程变得干净。
第三，边界条件分对称条件和表面对流条件两类。
